\begin{abstract}
Coding agents repeatedly reconstruct repository context, reason about changes,
and validate results. We study an obligation-centered supervision architecture
that retains source-linked semantic state, routes supported obligations through
checked symbolic mechanisms, and uses neural proposals for unresolved work.
Evidence applicability governs validation reuse and accepted-state publication.
The empirical question is whether their combined startup and operating costs
reduce neural work and validation effort while preserving useful progress.
We evaluate context substitution, avoided model calls, persistent planning,
validation reuse and native proof consumption as distinct mechanisms, retaining
an earlier additive-context cold-repair study as separate evidence.
On four constructed tasks, symbolic repair and residual routing pass the same
independent checks as a model-first baseline with 2,430 versus 4,865 provider
tokens. This 50.05\% reduction is confined to that development suite.
Other experiments exercise lifecycle-preserving reuse, memory consumption,
publication, crash recovery and proofs over actual saved state. Simple
retrieval matches memory's action reduction, and measured reuse savings do
not recover priming cost over five repetitions. In two development comparisons
on one natural two-issue repository sequence, every arm accepts one issue and
none completes the trajectory. These mixed findings establish bounded capabilities and
implementation limitations. Complete integrated token savings remain open;
component percentages cannot be combined into a system estimate.
\end{abstract}
